\uuid{nL9O}
\exo7id{1185}
\titre{exo7 1185}
\auteur{liousse}
\organisation{exo7}
\datecreate{2003-10-01}
\isIndication{false}
\isCorrection{false}
\chapitre{Déterminant, système linéaire}
\sousChapitre{Système linéaire, rang}
\cpge{mpsi-mp2i.calcul-matriciel-systemes.c, pcsi.calcul-matriciel-systemes-lineaires.c, ptsi.calcul-matriciel-systemes-lineaires.c, mpsi-mp2i.matrices-applications-lineaires.c, pcsi.matrices-applications-lineaires.d, ptsi.matrices-applications-lineaires.d, bcpst1.systemes-lineaires.a, bcpst1.systemes-lineaires.c, bcpst1.espace-vectoriel-kn.a}
\module{Algèbre}
\niveau{L2}
\difficulte{}

\contenu{
\texte{
On consid\`ere le syst\`eme $$(S) : \left \lbrace
\begin{array}{ll}
   x +  y+ z+  t &=0 \\
  x -  y-2z+2t &=0 \\
 2x + y+  z  &= 0 \\
\end{array} \right.$$
}
\begin{enumerate}
    \item \question{R\'esoudre le syst\`eme $(S)$ puis indiquer son rang.}
    \item \question{Montrer que l'ensemble des solutions de $(S)$ est un
sous-espace vectoriel de $\mathbb R^4$, indiquer sa dimension  
et en donner une base.}
\end{enumerate}
}
